
\chapter{מסקנות ועתיד החיזוי}

\section{ממצאים עיקריים}

מחקר זה הדגים כי \textenglish{xLSTM} \cite{beck2024xlstm} מצליח להתחרות ולעלות על ארכיטקטורות \textenglish{Transformer} מתקדמות בחיזוי טורי זמן רב-צעדי. שלושה ממצאים בולטים:

\begin{enumerate}
\item זיכרון מטריצה (\textenglish{mLSTM}) מאפשר לתפוס תלויות ארוכות טווח ביעילות שווה ל-\textenglish{attention} רב-ראשי.
\item מורכבות לינארית של \textenglish{xLSTM} יתרון משמעותי על רצפים ארוכים (\textenglish{n > 500}).
\item \textenglish{xLSTM} יציב יותר לשינויי \textenglish{learning rate} בהשוואה ל-\textenglish{Transformer}.
\end{enumerate}

\section{מגבלות ועבודה עתידית}

\textenglish{xLSTM} דורש אימון ממושך יותר עקב מורכבות המטריצה. כיוון מחקר עתידי הוא שילוב שכבות \textenglish{xLSTM} ו-\textenglish{attention} בארכיטקטורה היברידית, הניצולת את חוזקות שניהם: ייצוגיות מקומית (\textenglish{xLSTM}) ותפיסת תלויות גלובליות (\textenglish{Transformer}) \cite{nie2022patchtst}.
